\chapter{THỰC NGHIỆM VÀ THẢO LUẬN}

\section{Dữ liệu và cách chia tập}

Dữ liệu dùng trực tiếp trong notebook là gói reconia-data trên Kaggle được liên kết trong README của dự án, gồm SROIE, CORD-v2 và WildReceipt. SROIE có các trường company, address, date, total \parencite{huang2019}. CORD cung cấp chú thích văn bản, vị trí và nhóm ngữ nghĩa cho hóa đơn \parencite{park2019}. WildReceipt cung cấp nhãn cho các vùng chữ và được giới thiệu trong công trình SDMG-R \parencite{sun2021}.

\begin{table}[H]
\centering
\renewcommand{\arraystretch}{1.25}
\fontsize{12}{15}\selectfont
\setlength{\tabcolsep}{4pt}
\caption{Số mẫu dùng trong thực nghiệm}
\begin{tabularx}{\textwidth}{>{\raggedright\arraybackslash}X>{\centering\arraybackslash}X>{\centering\arraybackslash}X>{\centering\arraybackslash}X}
\toprule
\textbf{Bộ dữ liệu} & \textbf{Train} & \textbf{Validation} & \textbf{Test} \\
\midrule
SROIE & 501 & 125 & 347 \\
CORD-v2 & 800 & 100 & 100 \\
WildReceipt & 1014 & 253 & 472 \\
Tổng & 2315 & 478 & 919 \\
\bottomrule
\end{tabularx}
\end{table}

CORD-v2 giữ cách chia 800/100/100 có sẵn. SROIE và WildReceipt lấy 20\% tập train gốc làm validation với seed 42; giữ nguyên tập test. Tổng số ảnh test là 919. Không lấy mẫu test để huấn luyện hay chọn checkpoint.

Đối chiếu dữ liệu đã xác nhận CORD-v2 khớp bản phát hành của NAVER CLOVA và WildReceipt khớp gói OpenMMLab. Với SROIE, ảnh và text/box khớp bản Drive có bằng chứng lưu từ thời điểm phát hành; 347 GT test khớp từng byte với bản tham chiếu Kaggle sroie-datasetv2 version 4. Chưa đối chiếu trực tiếp được gói GT test từ cổng RRC vì không truy cập được. Đây là giới hạn xác minh nguồn, không phải kết luận rằng GT bị sai.


\section{Chuẩn bị dữ liệu đầu vào}
Với SROIE, nhóm chuyển các dòng chữ thành đơn vị từ và tạo nhãn cho bốn trường thông tin. Vị trí từ được suy ra từ vùng bao dòng nên không phải tọa độ từ được chú thích trực tiếp. Với CORD-v2, nhóm lấy văn bản, vị trí và loại thông tin từ chú thích có sẵn. Với WildReceipt, nhóm giữ các lớp giá trị cần trích xuất; các phần còn lại được gán nhãn O.

Cả ba mô hình dùng cùng dữ liệu đã chuyển đổi. Văn bản và vị trí đều lấy từ chú thích gốc; nhóm không chạy nhận dạng chữ lại trong các thí nghiệm này.

\section{Cấu hình huấn luyện}

Các lượt chạy dùng seed 42, AdamW, weight decay 0,01 cho nhóm tham số áp dụng decay, tối đa 30 epoch và early stopping với patience bằng 5. Effective batch size là 8: RoBERTa xử lý 8 mẫu mỗi batch; BROS và PICK xử lý 4 mẫu rồi tích lũy gradient qua 2 bước mới cập nhật trọng số. Checkpoint có F1 validation tốt nhất được lưu và nạp lại trước khi chấm test.

\begin{table}[H]
\centering
\renewcommand{\arraystretch}{1.25}
\fontsize{12}{15}\selectfont
\setlength{\tabcolsep}{4pt}
\caption{Cấu hình của các lượt huấn luyện đã hoàn thành}
\begin{tabularx}{\textwidth}{>{\raggedright\arraybackslash}X>{\centering\arraybackslash}X>{\centering\arraybackslash}X>{\raggedright\arraybackslash}X}
\toprule
\textbf{Mô hình} & \textbf{Batch} & \textbf{Tích lũy} & \textbf{Learning rate} \\
\midrule
RoBERTa & 8 & 1 & 3 \(\times\) \(10^{-5}\) \\
BROS & 4 & 2 & 3 \(\times\) \(10^{-5}\) \\
PICK & 4 & 2 & Nhóm backbone: \(10^{-4}\)\newline Nhóm còn lại: \(10^{-3}\) \\
\bottomrule
\end{tabularx}
\end{table}


Với PICK, nhóm tham số của nhánh ảnh được học với tốc độ nhỏ hơn các phần còn lại. Cách phân nhóm hiện tại còn đưa một lớp chiếu mới vào nhóm nhánh ảnh; ảnh hưởng của lựa chọn này chưa được đo riêng.

\section{Cách tính Precision Recall và F1}


Ba mô hình được chấm theo cùng cách: giá trị dự đoán phải khớp giá trị đúng và thuộc đúng trường sau khi chuyển chữ thường, gộp khoảng trắng. Mỗi giá trị dự đoán chỉ được ghép với một đáp án.

Precision cho biết trong các giá trị mô hình lấy ra, có bao nhiêu giá trị đúng. Recall cho biết trong các giá trị cần lấy, mô hình tìm được bao nhiêu. F1 tổng hợp hai mặt này. Gọi TP là số giá trị đúng, FP là số giá trị dự đoán sai và FN là số giá trị bị bỏ sót, các chỉ số của từng trường được tính bằng:
\begingroup
\setlength{\jot}{14pt}
\begin{align}
P &= \frac{\mathrm{TP}}{\mathrm{TP}+\mathrm{FP}}, \\
R &= \frac{\mathrm{TP}}{\mathrm{TP}+\mathrm{FN}}, \\
F_1 &= \frac{2\mathrm{TP}}{2\mathrm{TP}+\mathrm{FP}+\mathrm{FN}}.
\end{align}
\endgroup
Nếu mẫu số bằng 0, chỉ số tương ứng bằng 0. Nhóm tính từng chỉ số cho mỗi trường trên toàn bộ tập test rồi lấy trung bình giữa các trường. Do đó, F1 trung bình không nhất thiết bằng F1 tính lại từ Precision và Recall trung bình trong bảng.

Ví dụ, đáp án tổng tiền là ``21.60'' nhưng mô hình trả về ``21.6'' thì vẫn bị tính sai theo cách khớp chuỗi này. Kết quả dùng để so sánh ba mô hình trong đề tài; không đối chiếu trực tiếp với điểm công bố nếu bài gốc dùng cách chấm khác. Riêng WildReceipt, nhóm chấm giá trị trích xuất, khác cách chấm nhãn từng vùng của tác vụ gốc.

\section{Kết quả trên tập test}

\begin{table}[H]
\centering
\renewcommand{\arraystretch}{1.25}
\fontsize{12}{15}\selectfont
\setlength{\tabcolsep}{4pt}
\caption{Precision Recall và F1 trên tập test với seed 42 (làm tròn bốn chữ số thập phân)}
\begin{tabularx}{\textwidth}{>{\raggedright\arraybackslash}X>{\raggedright\arraybackslash}X>{\centering\arraybackslash}X>{\centering\arraybackslash}X>{\centering\arraybackslash}X}
\toprule
\textbf{Bộ dữ liệu} & \textbf{Mô hình} & \textbf{Precision} & \textbf{Recall} & \textbf{F1} \\
\midrule
SROIE & RoBERTa & 0,8682 & 0,8710 & 0,8696 \\
\textbf{SROIE} & \textbf{BROS} & \textbf{0,8848} & \textbf{0,8811} & \textbf{0,8829} \\
SROIE & PICK & 0,5568 & 0,6001 & 0,5619 \\
CORD-v2 & RoBERTa & 0,7651 & 0,8039 & 0,7761 \\
\textbf{CORD-v2} & \textbf{BROS} & \textbf{0,7918} & \textbf{0,8221} & \textbf{0,7982} \\
CORD-v2 & PICK & 0,3901 & 0,3783 & 0,3731 \\
WildReceipt & RoBERTa & 0,8684 & 0,8294 & 0,8319 \\
\textbf{WildReceipt} & \textbf{BROS} & \textbf{0,9238} & \textbf{0,9167} & \textbf{0,9185} \\
WildReceipt & PICK & 0,6290 & 0,6450 & 0,6308 \\
\bottomrule
\end{tabularx}
\end{table}

BROS có F1 cao nhất trong cả ba bộ, vượt RoBERTa lần lượt 1,33; 2,21 và 8,66 điểm phần trăm. PICK có F1 thấp hơn hai mô hình còn lại trong các cấu hình đã chạy. 

\begin{table}[H]
\centering
\renewcommand{\arraystretch}{1.25}
\fontsize{12}{15}\selectfont
\setlength{\tabcolsep}{4pt}
\caption{Epoch đã chạy và epoch được chọn theo validation}
\begin{tabularx}{\textwidth}{>{\raggedright\arraybackslash}X>{\centering\arraybackslash}X>{\centering\arraybackslash}X>{\centering\arraybackslash}X}
\toprule
\textbf{Bộ dữ liệu} & \textbf{RoBERTa} & \textbf{BROS} & \textbf{PICK} \\
\midrule
SROIE & 14 / 9 & 17 / 12 & 27 / 22 \\
CORD-v2 & 17 / 12 & 28 / 23 & 24 / 19 \\
WildReceipt & 24 / 19 & 13 / 8 & 20 / 15 \\
\bottomrule
\end{tabularx}
\end{table}

Mỗi ô ghi tổng epoch đã chạy / epoch tốt nhất. Cả chín lượt dừng sau năm epoch không cải thiện, phù hợp với patience đã đặt. Điều này xác nhận hành vi dừng sớm trong log, nhưng không chứng minh từng mô hình đã được tối ưu đầy đủ.

\section{Ví dụ phân tích dự đoán sai}

Nhóm xuất lại dự đoán từ chín checkpoint trên Kaggle để có kết quả cho từng ảnh và từng đoạn chữ. Có 2.757 bản ghi dự đoán theo ảnh, bằng 919 ảnh test nhân ba mô hình. Lượt xuất bổ sung không thay đổi dự đoán trường--giá trị đã lưu; chỉ bổ sung nhãn của từng đoạn và trạng thái đã được mô hình xử lý hay chưa.

\begin{table}[H]
\centering
\renewcommand{\arraystretch}{1.25}
\fontsize{12}{15}\selectfont
\setlength{\tabcolsep}{4pt}
\caption{Một số lỗi của RoBERTa trên mẫu test CORD-v2 có ID 0}
\begin{tabularx}{\textwidth}{>{\raggedright\arraybackslash}X>{\raggedright\arraybackslash}X>{\raggedright\arraybackslash}X}
\toprule
\textbf{Văn bản} & \textbf{Nhãn GT} & \textbf{Nhãn dự đoán} \\
\midrule
60.000 & B-menu.price & B-menu.unitprice \\
60.000 & B-menu.itemsubtotal & B-menu.price \\
-60.000 & B-sub\_total.discount\_price & B-sub\_total.etc \\
xx7730 & O & B-total.creditcardprice \\
\bottomrule
\end{tabularx}
\end{table}

Ví dụ cho thấy mô hình có thể giữ đúng chữ nhưng gán sai ý nghĩa. Hai đoạn “60.000” là hai vị trí khác nhau trong cùng hóa đơn, mang hai nhãn GT khác nhau. Vì thế, chỉ nhìn danh sách giá trị tiền khó nhận ra nguyên nhân; cần đối chiếu vị trí và nhãn từng đoạn. Một mẫu minh họa không đủ để đại diện cho toàn bộ tập test hoặc chứng minh nguyên nhân chênh lệch giữa các kiến trúc.


\section{Thảo luận và giới hạn}
BROS đạt F1 cao nhất trên cả ba bộ dữ liệu trong các lượt đã chạy. Việc dùng thêm vị trí có thể hỗ trợ phân biệt các trường, nhưng nhóm chưa có thí nghiệm tách riêng yếu tố này. Kết quả chưa đủ để kết luận BROS luôn tốt hơn trong mọi loại hóa đơn.

PICK có điểm thấp hơn hai mô hình còn lại. Do bản triển khai khác bài gốc ở một số bước, nhóm cần kiểm tra lại trước khi xác định nguyên nhân. Kết quả hiện tại chưa đủ để kết luận hướng dùng ảnh và đồ thị kém hiệu quả.

Các thí nghiệm mới dùng một giá trị khởi tạo ngẫu nhiên là 42, chưa đánh giá độ dao động qua nhiều lần chạy. Giới hạn độ dài làm RoBERTa cắt 2 ảnh và BROS cắt 1 ảnh trên SROIE; trên WildReceipt, mỗi mô hình cắt 1 ảnh. PICK không có mẫu vượt giới hạn đã đặt. Ảnh hưởng của việc cắt này chưa được đo riêng.

So sánh chính là giữa các mô hình trên cùng bộ dữ liệu. Không dùng điểm cao hơn ở một bộ để khẳng định bộ đó luôn dễ hơn, vì số trường và cách chú thích giữa các bộ khác nhau.
